\section{Related work}
\label{sec:related}

% STE descriptive
Table~\ref{tab:methods} compares the seven studies closest to this one. It gives the task, the
algebra, the moduli, the coordinates of the analysis and the kind of evidence that each study
reports.

\begin{table}[t]
  \centering
  \footnotesize
  \caption{\textbf{What each method can see.} Each row is one study, and every cell comes
    from that study's PDF. The evidence column gives the tests that the study reports. A null
    arm is a control condition for which the account predicts no effect. ``Not stated'' means
    that the PDF does not contain the word pre-registration.}
  \label{tab:methods}
  \setlength{\tabcolsep}{3pt}
  \renewcommand{\arraystretch}{1.15}
  \begin{tabular}{L{0.10\textwidth} L{0.055\textwidth} L{0.12\textwidth} L{0.14\textwidth}
      L{0.13\textwidth} L{0.15\textwidth} L{0.12\textwidth} L{0.075\textwidth}}
    \toprule
    study & task & algebra & moduli (square-free?) & coordinates & evidence & null arm & pre-reg. \\
    \midrule
    \citet{nanda2023progress} & $a + b$ & $\mathbb{Z}/p\mathbb{Z}$ under $+$ & $113$, $53$, $109$ (primes) & Fourier basis & logit ablation; activation replacement & non-key frequencies & not stated \\
    \citet{zhong2023clock} & $a + b$ & $\mathbb{Z}/p\mathbb{Z}$ under $+$ & $59$ (prime) & Fourier basis; principal components & correlational; rank-2 embedding replacement & none reported & not stated \\
    \citet{chughtai2023toy} & $a b$ in a group & finite groups & $C_{113}$, $C_{118}$ (square-free), and five non-cyclic groups & irreducible representations & logit ablation; activation replacement & directions predicted to be noise & not stated \\
    \citet{nguyen2026dlogclock} & $a \cdot b$ & $(\mathbb{Z}/p\mathbb{Z})^{*}$ & $113$; ten primes from $59$ to $113$ & discrete-log characters & correlational & none reported & not stated \\
    \citet{chen2026multiplication} & $a \cdot b$ & the monoid $(\mathbb{Z}/n\mathbb{Z}, \cdot)$ & $113$, $143$, $154$, $165$ (all square-free) & characters of each $\mathcal{J}$-class group & correlational, in their words & random vocabulary subsets & not stated \\
    \citet{hwang2026pfe} & $a + b$ & $\mathbb{Z}/N\mathbb{Z}$ under $+$, by CRT & primes; ten composites from $15$ to $385$ (all square-free) & prime-indexed Fourier features, fixed & input-channel ablation, MLP classifier & irrelevant channels; shuffled features & not stated \\
    this work & $a \cdot b$ & the monoid $(\mathbb{Z}/n\mathbb{Z}, \cdot)$ & $23$ moduli, $14$ not square-free & characters of each stratum's local group $G_{m}$ & logit ablation; weight intervention at $n = 113$, $121$ & random character sets & $9$ of the $14$ rows of Table~\ref{tab:glance} \\
    \bottomrule
  \end{tabular}
\end{table}
% STE end

\paragraph{Grokking as an instrument.} \citet{power2022grokking} established the phenomenon:
on small algorithmic datasets, generalisation can appear long after the training set is fit,
and the smaller the training fraction the longer the delay. That paper claims a testbed rather
than a mechanism, and the testbed is what the literature since has used.
\citet{nanda2023progress} supplied the mechanism for $a + b \bmod n$ by reverse-engineering a
one-layer transformer completely: the network embeds inputs in Fourier features, adds phases
through the attention and MLP, and reads out with a trigonometric identity, confirmed by
ablations in Fourier space. Their restricted and excluded losses are the ablation protocol we
use, in a coordinate they did not need.

\paragraph{The algorithm is not inevitable.} \citet{zhong2023clock} showed that the Clock
algorithm is not the only one a network may find on the same task: a Pizza algorithm is
prevalent under other hyperparameters, algorithmic phase transitions exist, and networks
sometimes ensemble. \citet{chughtai2023toy} generalised the setting to arbitrary finite groups
and found universality to be two-level: the circuit \emph{family} is characterisable by
representation theory, while the particular circuit and the order in which it develops are
arbitrary. Their paired-ablation design, testing both the signal the theory predicts and the
noise it predicts, is the standard we hold our own ablations to.
\citet{stander2024cosets} analyse the same territory through cosets.

\paragraph{Multiplication and the basis.} Multiplication modulo $n$ resisted this treatment
for a simple reason: in raw residue coordinates a grokked multiplicative model has a dense
spectrum. That observation is \citet{furuta2024towards}'s, whose Section~5.2 reports that
multiplication ``leverages all frequencies'' with cosine-biased components, against the sparse
key-frequency picture addition gives.
\citet{nguyen2026dlogclock} showed the density is a basis artefact. Re-analysed
in the multiplicative character basis, where discrete log turns multiplication into addition,
the same model is sparse, and they report $\mathrm{Gini}\ 0.58$ against $0.07$, four key
frequencies, and $96.9\%$ of neurons single-frequency. Section~\ref{sec:calib} replicates
their spectral statistics and withdraws one clause of our own earlier replication.
\citet{doshi2024grokking} classify which modular polynomials are learnable and offer it as a
hypothesis with supporting evidence; their inverse participation ratio is a third convention,
reciprocal and bounded in $[0,1]$, which is why Section~\ref{sec:methods-sparsity} states ours
explicitly. \citet{furuta2024towards} and \citet{kunin2024getrich} supply the initialisation
and lazy-to-rich picture our initialisation control follows.
The lesson generalises past spectra, and we state it because the cost is easy to underestimate:
\emph{any} diagnostic calibrated on addition can read a multiplicative model as structureless
while the structure is present in exponent coordinates. We hit this ourselves. The same
grokked model's top eight logit components read $17.0\%$ ``$a+b$'' raw and $98.4\%$ in discrete-log coordinates
(\S\ref{sec:methods-sparsity}), and we record it as a methodological hazard rather than a
finding about multiplication.

\paragraph{Divisor structure is something these models already organise around.}
Before any character account, \citet{charton2024learning} trained transformers to compute
$\gcd(a, b)$ and characterised their predictions completely. The model learns a set $D$ of
integers and returns the largest element of $D$ dividing both inputs, so that every pair
sharing a $\gcd$ receives the same output, \emph{including where that output is wrong}.
The equivalence structure therefore sits in the computation rather than only in the target.
That is the same shape as Consequence~2 of Theorem~\ref{thm:stratum}, reached on a different
task, in base-$10$ digit coordinates, and with no representation-theoretic model; our fibre
residual \eqref{eq:resid} is the quantitative version of the constancy he observes
qualitatively. We read it as independent evidence that divisor classes are a structure these
models find, and note that he also reports small primes being grokked only after extended
training.

\paragraph{The stratified framing, which is not ours.}
\citet{chen2026multiplication} extend the character account past the group case. Where global
invertibility fails they model $(\mathbb{Z}/n\mathbb{Z}, \cdot)$ as a monoid, show that
group-character-style computation \emph{localises to disjoint algebraic strata that partition
the input space}, identify those strata with Green's $\mathcal{J}$-classes, and demonstrate
class-sensitive attention routing. They do so on four moduli, $113$, $143$, $154$ and $165$,
every one of which their Table~1 marks square-free. Their Theorem~3.4 requires each class to \emph{be} a group,
\begin{equation}
  \label{eq:chen}
  J_{d} \;\cong\; (\mathbb{Z}/(n/d)\mathbb{Z})^{\times}
  \quad\text{with a local identity } e_{J_{d}} \text{ and a local inverse } c^{\sharp},
\end{equation}
and their account of the read-out, developed in their Section~4 rather than in
Theorem~3.4, evaluates $\mathrm{Logit}(c) \propto \chi_{\rho_{d}}(a\,b\,c^{\sharp})$ on that
group. A non-regular class supplies neither: with no idempotent there is no $c^{\sharp}$.
Proposition~\ref{prop:torsor} replaces the isomorphism in \eqref{eq:chen} by a
\emph{bijection} that needs no group structure, and Theorem~\ref{thm:stratum} recovers the
character account on $G_{m}$ instead of on $J_{d}$. This is not a counterexample to Clifford--Munn--Ponizovski\u{\i}. CMP
\citep{steinberg2016monoids}, which \citet{chen2026multiplication} state as their
Theorem~D.16, says that the irreducible representations of a finite monoid are indexed by the
\emph{regular} $\mathcal{J}$-classes: a non-regular class contributes none, and no argument
here disputes that. What Theorem~\ref{thm:stratum} factors through $G_{m}$ is the
\emph{function the network has to compute} on a stratum, not the monoid's representation
theory. The characters we measure are characters of the local group $G_{m}$, which is a
genuine group; they are not irreps of $(\mathbb{Z}/n\mathbb{Z}, \cdot)$ attached to a
non-regular class, and the reduction does not require any. What is new here is not a stronger theorem but a wider one, and the width is the
contribution. Proposition~\ref{prop:torsor} trades a group isomorphism for a bijection of
sets: weaker hypotheses and a weaker conclusion, which is a different statement and not a
strengthening of their Theorem~3.4, and we do not claim it as one. Two things follow that
their theorem does not give. \emph{First, non-vacuity exactly where theirs is vacuous.}
\eqref{eq:chen} asks $J_{d}$ to be a group, so at a non-square-free modulus it is silent on
precisely the classes that make such a modulus interesting: $8$ of the $16$ at $n = 120$,
and the ones containing the nilpotents. \eqref{eq:torsor} and Theorem~\ref{thm:stratum}
are stated for every $d \mid n$ and ask nothing beyond $\mathcal{J}$-classhood. \emph{Second, cross-stratum composition.}
Theorem~\ref{thm:stratum} is a statement about the pair $J_{d} \times J_{e}$ for all
$d, e$, which their Theorem~3.4 does not address even in the square-free case. It says which
class a block lands in, \eqref{eq:jcompose}, and which local group carries the
multiplication once it is there, \eqref{eq:stratum}. That is the statement our
measurements use, because the strata of the input grid are \emph{blocks} and not
classes, and every per-stratum number in Section~\ref{sec:strata} is measured on a block.
The framing is theirs, and this
paper is its completion, not its rediscovery. They state two limitations of their own result in their own
words: that their evidence is only correlational, and that non-square-free moduli introduce
non-regular $\mathcal{J}$-classes containing nilpotent elements, which they describe as
breaking the reduction to local group characters. Sections~\ref{sec:causal} and
\ref{sec:strata} supply the two.

The relation to \citet{chughtai2023toy} is worth naming precisely, because it is the shape of
our result. They argue universality is two-level: the circuit family invariant, the specific
circuit arbitrary. What we add is that for $a \cdot b \bmod n$ the two levels are
\emph{algebraic}: the family is the character clock, and the index set is fixed by the
stratum's own local group.

\paragraph{The CRT connection is not ours either.} \citet{hwang2026pfe} supply a
prime-indexed embedding and prove block-diagonality by Schur's lemma, with an ablation carrying
a null arm: task-relevant channels drop $0.60$--$0.92$ while irrelevant channels sit at or
below the noise floor. Theirs is the basis \emph{supplied}. What Section~\ref{sec:crt-law}
adds is that the equivalent structure is \emph{learned} rather than imposed, on
\emph{multiplication} rather than addition, and at \emph{non-square-free} moduli, where the
predicted set is built from maximal prime powers rather than primes and the distinction is
testable.
